\subsection{INFINITE PRODUCTS}

To each sequence \((\pmb{x}_n)\) of points in a normed algebra A, let us make correspond the sequence of partial products \(p_n = \prod_{k=0}^{n} x_k\) ; then the pair of sequences \((x_{n})\) and \((p_{n})\) is called the infinite product whose general factor is \(x_{n}\) . The infinite product with general factor \(x_{n}\) is said to be convergent if the sequence \((p_{n})\) is convergent in A; the limit of this sequence is then called the product of the sequence \((x_{n})\) and is denoted by \(\bigoplus_{n=0}^{\infty} x_{n}\) .

PROPOSITION 2. Let \((\pmb{x}_n)\) be a sequence of points in a complete normed algebra A. a) If the infinite product whose general factor is \(\pmb{x}_n\) is convergent and if \(\mathbf{P}_{n=0}^{\infty}\pmb{x}_n\) is a unit in A, then for each \(\varepsilon > 0\) there exists an integer \(n_0\) such that

\[
\left\| \prod_ {k = m} ^ {n} x _ {k} - e \right\| \leqslant \varepsilon
\]

whenever \(n_0 \leqslant m \leqslant n\) .

\begin{enumerate}
\def\labelenumi{\alph{enumi})}
\setcounter{enumi}{1}
\tightlist
\item
  Conversely, if the sequence \((\pmb{x}_n)\) satisfies this condition, the infinite product with general factor \(\pmb{x}_n\) is convergent; and if each of the \(\pmb{x}_n\) is a unit in A, then \(\mathbf{P}_{n=0}^{\infty}\pmb{x}_n\) is a unit.
\end{enumerate}

The proof of this proposition follows step by step the proof of Theorem 1, and is left to the reader (the finite subsets L of N in the proof of Theorem 1 are to be replaced by intervals).

COROLLARY I. If the infinite product with general factor \(x_{n}\) is convergent, and if \(\bigcap_{n=0}^{\infty}x_{n}\) is a unit, then \(\lim_{n\to\infty}x_{n}=e\) .

COROLLARY 2. If the infinite product with general factor \(\pmb{x}_n\) is convergent, and if \(\mathop{\sum }\limits_{{n = 0}}^{\infty}{\pmb{x}}_{n}\) is a unit, then the infinite product with general factor

\[
\mathbf {y} _ {n} = \mathbf {x} _ {n + h} \quad (n \geqslant 0)
\]

is convergent.

The product of the sequence \((\mathbf{y}_n)\) is denoted by \(\bigoplus_{n = h}^{\infty}\mathbf{x}_n\) , and is also called the residue of index \(h\) of the infinite product with general factor \(\mathbf{x}_n\) .

Still under the assumption that \(\mathbf{P}_{n=h} x_n\) is a unit, it follows from Proposition 2 that if \((z_n)\) is a sequence such that \(z_n = x_n\) for all but a finite number of indices, then the product with general factor \(z_n\) is convergent.

PROPOSITION 3. Let \((k_{n})\) be a strictly increasing sequence of integers \(\geqslant 0\) , such that \(k_{0} = 0\) ; if the infinite product with general factor \(x_{n}\) converges, and

if we put

\[
\boldsymbol {u} _ {n} = \prod_ {p = k _ {n}} ^ {k _ {n + 1} - 1} \boldsymbol {x} _ {p},
\]

then the infinite product whose general factor is \(u_{n}\) is convergent and we have

\[
\underset {n = 0} {\overset {\infty} {\mathsf {P}}} u _ {n} = \underset {n = 0} {\overset {\infty} {\mathsf {P}}} x _ {n}.
\]

For the sequence of partial products of the sequence \((\boldsymbol{u}_{n})\) is a subsequence of the sequence of partial products of the sequence \((\boldsymbol{x}_{n})\) .

Finally, the same argument as was used for abelian groups (Chapter III, § 5, no. 7) shows that if a sequence \((\pmb{x}_n)\) in a normed algebra \(\mathbf{A}\) is multipleiable, then the product whose general factor is \(\pmb{x}_n\) is convergent, and

\[
\mathop {\text { P }} _ {n = 0} ^ {\infty} x _ {n} = \prod_ {n \in \mathbf {N}} x _ {n}
\]

(which is also written as \(\prod_{n=0}^{\infty} x_n\) ); the converse is of course not true (cf.~Exercise 7).
% semantic-fragments: legacy-input-seam
\relax{}
\ExerciseGroupsStart
